\scanpage{360}
\mspage{99}

10. \uline{Questions de comparaison \unsure{pour} les catégories Cat.-\unsure{fibrée|fibrées} \deleted{\unsure{Cat}} satisfaisant au ``th. d'homotopie''} (base \unsure{filtrante ?|filtrante}).

Nous allons généraliser maintenant au cas d'une catégorie (Cat)-fibrée sur une autre
\diagram{\begin{tikzcd}[ampersand replacement=\&]
C\arrow[d,"\unsure{\pi|p}"']\\ T,
\end{tikzcd}}
\margin{\unsure{(}}
en pensant surtout au cas $T=\underline{\mathrm{Ty}}$, $C=\underline{\mathrm{Ch}}_\Sigma(X;a,b)$. L'hypothèse essentielle est la suivante

\margin{(Hypothèse d'homotopie)}Si $t,t'\in\mathrm{Ob}\,T$ sont tels que $\mathrm{Hom}(t,t')\ne\varnothing$, alors les foncteurs $\varphi^*:C_{t'}\mathrel{\unsure{-|\to}}C_t$ associés aux $\varphi\in\mathrm{Hom}(t,t')$ sont homotopes.

\deleted{\unsure{Soit donc}} Ils définissent donc une seule et \unsure{$=$|même} flèche dans $\mathrm{Cat\,hot}$. \deleted{\unsure{So}} Je désigne par $T_{\mathrm{ord}}$ \inserted{ou $T_0$} l'ens $\mathrm{Ob}(T)$, avec comme relation d'ordre celle déduite de la structure de catégorie :
\[
t\leq t'\quad\Longleftrightarrow\quad\mathrm{Hom}_T(t,t')\ne\varnothing,
\]
\deleted{\unsure{donc (ainsi} $T_{\mathrm{ord}}$} de sorte qu'on a un foncteur
\[
T\longrightarrow T_0
\]
qui est \uline{universel} pour les foncteurs de $T$ dans des catégories ``ordonnées'' (i.e. où \unsure{\deleted{deux}|toute} flèche \unsure{est} mono-morphique) \inserted{(ou pour les foncteurs $F$ sur $T$ tels que $F$ soit constant sur tout $\mathrm{Hom}(t,t')$)}. L'\unsure{hypth.|hypothèse|hyp.} d'homotopie peut s'exprimer en disant que le \deleted{\unsure{a lieu de la manière toute}} \inserted{foncteur composé}
\diagram{\begin{tikzcd}[ampersand replacement=\&,column sep=small,row sep=small]
\deleted{\unsure{C}}\,T^\circ\arrow[r]\&\mathrm{Cat}\arrow[r]\&\mathrm{Cat\,hot}\\
t\arrow[r,mapsto]\& C_t\&
\end{tikzcd}}

% archive sheet number: 359
